\section{Theoretical Analysis}
\label{app:theoretical_analysis}

Here, we provide a detailed analysis of the empirical results obtained in the main text from a theoretical perspective. We begin by repeating the derivation of the baseline compute-optimal scaling for the number of tokens as a function of model parameters, and continue to systematically work through the perturbations discussed in~\cref{sec:robustness_analysis}.

\subsection{Baseline Compute-Optimal Scaling Derivation}
\label{app:theoretical_analysis:subsec:baseline}

The Chinchilla scaling law for pretraining loss $L$ as a function of non-embedding model parameters $N$ and number of training tokens $D$ is given by~\cref{eqn:scaling_law}, which we repeat here
\begin{equation}
    L(N, D) = E + \frac{A}{N^\alpha} + \frac{B}{D^\beta},
\end{equation}
where $E$ is the irreducible loss, and $(A, \alpha)$ and $(B, \beta)$ are parameters for the model size and data scaling terms, respectively.

The training compute budget $C$ is approximately proportional to the product of model size and training data, $C \approx c ND$, where $c>0$ is some constant factor. This allows us to express the number of training tokens as a function of compute and model size: $D = C / (c N)$. Substituting this into the loss function yields the loss for a fixed compute budget
\begin{equation}
\label{eq:fixed_budget_loss}
    L(N, C) = E + A N^{-\alpha} + B \left(\frac{C}{c N}\right)^{-\beta} = E + A N^{-\alpha} + B(c^\beta C^{-\beta}) N^\beta.
\end{equation}

The optimal model size $N_{\text{opt}}$ that minimizes the loss for a fixed compute budget $C$ is simply found by differentiating~\cref{eq:fixed_budget_loss} with respect to $N$ and setting the derivative to zero
\begin{equation}
    \frac{\partial L}{\partial N} = - \alpha A N^{-(\alpha+1)} + \beta B (c^\beta C^{-\beta}) N^{\beta-1} = 0.
\end{equation}
Rearranging this equation reveals the optimal trade-off:
\begin{equation}
    \alpha A N_{\text{opt}}^{-(\alpha+1)} = \beta B (c^\beta C^{-\beta}) N_{\text{opt}}^{\beta-1}.
\end{equation}
Solving for $N_{\text{opt}}$:
\begin{equation}
    N_{\text{opt}}^{\alpha+\beta} = 
    \frac{\alpha A}{\beta B c^\beta} C^\beta 
    \implies 
    N_{\text{opt}} =
    \left( \frac{\alpha A}{\beta B c^\beta} \right)^{\frac{1}{\alpha+\beta}} C^{\frac{\beta}{\alpha+\beta}}
\end{equation}

The compute-optimal tokens-per-parameter ratio is $R_{\text{opt}} = D_{\text{opt}} / N_{\text{opt}}$. Using $D_{\text{opt}} = C / (c N_{\text{opt}})$, we find $R_{\text{opt}} = C / (c N_{\text{opt}}^2)$. Substituting our expression for $N_{\text{opt}}$ results in 
\begin{equation}
    R_{\text{opt}} = \frac{C}{c} \left[ \left( \frac{\alpha A}{\beta B c^\beta} \right)^{\frac{1}{\alpha+\beta}} C^{\frac{\beta}{\alpha+\beta}} \right]^{-2} = \frac{C}{c} \left( \frac{\beta B c^\beta}{\alpha A} \right)^{\frac{2}{\alpha+\beta}} C^{\frac{-2\beta}{\alpha+\beta}}.
\end{equation}
This simplifies to the final form, which shows the ratio's dependence on the compute budget $C$ as
\begin{equation}
    \frac{D_{\text{opt}}}{N_{\text{opt}}} = K \cdot C^{\frac{\alpha-\beta}{\alpha+\beta}},
    \label{eq:ratio}
\end{equation}
where $K = \frac{1}{c} \left( \frac{\beta B c^\beta}{\alpha A} \right)^{\frac{2}{\alpha+\beta}}$ is a constant. The key insight from the Chinchilla works is that empirically, one finds that $\alpha \approx \beta$, making the exponent on $C$ approximately zero and the optimal ratio nearly constant. In mathematical terms, it leads to the relation
\begin{equation}
    \frac{D_{\text{opt}}}{N_{\text{opt}}} \approx K =
    % \frac{1}{6}
    % \left( \frac{ B 6^\alpha}{ A} \right)^{\frac{1}{\alpha}}.
    \left( \frac{ B }{ A} \right)^{\frac{1}{\alpha}}.
    \label{eq:empirical_ratio}
\end{equation}
If we want to recover the $20:1$ ratio of training tokens to number of parameters we expect to find that $B \approx 2.85 A$ for $\alpha \approx 0.35$.

\subsection{Analysis of Perturbations}
We now analyze how systematic errors in model parameter counts affect the fitted scaling parameters $(\hat{A}, \hat{\alpha})$ and the resulting optimal ratio. Let $N$ be the true parameter count and $\tilde{N}$ be the perturbed (incorrect) count used for fitting. The fitting process minimizes the error between the model $L(\tilde{N}, D) = \hat{E} + \hat{A}\tilde{N}^{-\hat{\alpha}} + \hat{B}D^{-\hat{\beta}}$ and the observed losses. For most types of perturbations we consider, since the data $D$ is unaffected, we assume $\hat{B} \approx B$ and $\hat{\beta} \approx \beta$. We will break this assumption when necessary.

\subsubsection{Multiplicative Constant Perturbation}
\label{app:theoretical_analysis:subsec:multiplicative}

Assume the reported parameters are a constant multiple of the true parameters: $\tilde{N} = c_m \cdot N$. The model size term in the loss is modified to
\begin{equation}
\label{eqn:multi_constant}
    \hat{A} \tilde{N}^{-\hat{\alpha}} = \hat{A} (c_m N)^{-\hat{\alpha}} = (\hat{A} c_m^{-\hat{\alpha}}) N^{-\hat{\alpha}}.
\end{equation}
For \cref{eqn:multi_constant} to match the true term $A N^{-\alpha}$, the fitting procedure will ideally find parameters such that:
\begin{itemize}
    \item $\hat{\alpha} \approx \alpha$
    \item $\hat{A} c_m^{-\hat{\alpha}} \approx A \implies \hat{A} \approx A c_m^{\alpha}$
\end{itemize}
The exponent in the optimal ratio (\ref{eq:ratio}) becomes $(\hat{\alpha} - \beta)/(\hat{\alpha} + \beta) \approx (\alpha - \beta)/(\alpha + \beta)$.

\textbf{Conclusion:} A multiplicative error does not change the exponent governing the trend of the optimal ratio. The flat relationship with compute budget is preserved. However, the constant prefactor $K$ is shifted by a factor of $(c_m^\alpha)^{-2/(\alpha+\beta)} = c_m^{-2\alpha/(\alpha+\beta)}\approx c_m^{-1}$, which shifts the entire line up or down on a log-log plot. This is shown in Fig.~\ref{fig:perturbed_fits} (top row) and Fig.~\ref{fig:perturbed_compute_optimal_tpp} (top left).

%%%START THEO
%\theo{Alternative proof - slightly more rigorous}

%Consider Chinchilla's paper equation 11 (also find below for completeness).

%Let 
%\begin{align}
%F(a, b, \alpha, \beta, e) &= \sum_{i} \text{Huber}_\delta(\text{LSE}_e(a - \alpha \log N_i, b - \beta \log D_i) - \log L_i) \\
%F_{c_m}(a, b, \alpha, \beta, e) &= \sum_{i} \text{Huber}_\delta(\text{LSE}_e(a - \alpha \log (c_m\cdot N_i), b - \beta \log D_i) - \log L_i)
%\end{align}

%If $(a^*, b^*, \alpha^*, \beta^*, e^*)$ minimizes $F$, then for any point $(a, b, \alpha, \beta, e)$:
%\begin{equation}
%F(a^*, b^*, \alpha^*, \beta^*, e^*) \leq F(a, b, \alpha, \beta, e)
%\end{equation}


%Knowing that $F_{c_m}(a, b, \alpha, \beta, e) = F(a - \alpha \log(c_m), b, \alpha, \beta, e)$, for any $(a, b, \alpha, \beta, e)$:
%\begin{align}
%F_{c_m}(a^* + \alpha^* \log(c_m), b^*, \alpha^*, \beta^*, e^*) &= F(a^*, b^*, \alpha^*, \beta^*, e^*) \\
%&\leq F(a - \alpha \log(c_m), b, \alpha, \beta, e) \\
%&= F_{c_m}(a, b, \alpha, \beta, e)
%\end{align}

%Therefore $(a^* + \alpha^* \log(c_m), b^*, \alpha^*, \beta^*, e^*)$ minimizes $F_{c_m}$.

%Conversely, if $(\tilde{a}, \tilde{b}, \tilde{\alpha}, \tilde{\beta}, \tilde{e})$ minimizes $F_{c_m}$, then $(\tilde{a} - \tilde{\alpha} \log(c_m), \tilde{b}, \tilde{\alpha}, \tilde{\beta}, \tilde{e})$ minimizes $F$.

%This establishes a bijection between minimizers (assuming their existence). Using the re-parametrization $A=e^{a}$, we conclude that the new solution to the (perturbed) optimization problem is $e^{(a + \alpha \log cm)} = A c_m^{\alpha}$
%%%END THEO




\subsubsection{Additive Constant Perturbation}
\label{app:theoretical_analysis:subsec:additive}

Assume an additive error, e.g., from including/excluding embeddings: $\tilde{N} = N + c_a$. The loss term $\hat{A}(N+c_a)^{-\hat{\alpha}}$ is no longer a pure power law in $N$.

We examine the process of fitting the perturbed model $g(N; \hat{A}, \hat{\alpha}) = \hat{A}(N+c_a)^{-\hat{\alpha}}$ to the true function $f(N) = AN^{-\alpha}$ by minimizing the Mean Squared Error (MSE) in log-space.

The objective is to find $(\hat{A}, \hat{\alpha})$ that minimize $\sum_i [ \log f(N_i) - \log g(N_i) ]^2$.
The core of the problem lies in approximating the term $\log(N+c_a)$. For the regime where $N \gg |c_a|$, which applies to the larger models in the study, we can analyze the local behavior of the function.

\paragraph{Effect on the Scaling Exponent $\hat{\alpha}$:}
The most critical parameter in a power law is its exponent, which corresponds to the slope in a log-log plot. The true slope is constant
\begin{equation}
    \frac{d(\log f(N))}{d(\log N)} = -\alpha
\end{equation}
For the perturbed function, the effective slope is not constant and depends on $N$ as
\begin{align}
    \text{Effective Slope}(N) &= \frac{d(\log g(N))}{d(\log N)} = -\hat{\alpha} \frac{d(\log(N+c_a))}{d(\log N)} = -\hat{\alpha} \left( \frac{N}{N+c_a} \right)
    \\ \nonumber
    &\implies \hat{\alpha}=\alpha/\left( \frac{N}{N+c_a} \right).
\end{align}
The fitting procedure must select a single exponent $\hat{\alpha}$ that best represents this varying slope over the range of data. To match the true average slope of $-\alpha$, the fitted $\hat{\alpha}$ must compensate for the factor $N/(N+c_a)$:
\begin{itemize}
    \item When $c_a > 0$, the factor $N/(N+c_a) < 1$. To achieve the target slope, the fitting process must select an exponent $\hat{\alpha} > \alpha$.
    \item When $c_a < 0$, the factor $N/(N+c_a) > 1$ (for $N>|c_a|$). To compensate, the fitting process must select an exponent $\hat{\alpha} < \alpha$.
\end{itemize}
This provides a direct analytical explanation for the observed behavior of $\hat{\alpha}$ in Figure 4, which is smaller than the true $\alpha$ for $c_a<0$ and increases approximately linearly for $c_a>0$.

\paragraph{Effect on the Prefactor $\hat{A}$:}
Once the optimal $\hat{\alpha}$ is determined, the prefactor $\hat{A}$ is chosen to minimize the remaining offset. We can approximate this by enforcing that the functions match at some effective "pivot" point $N_0$ that is characteristic of the dataset.
\begin{equation}
    f(N_0) \approx g(N_0) \implies A N_0^{-\alpha} \approx \hat{A} (N_0+c_a)^{-\hat{\alpha}}.
\end{equation}
Solving for $\hat{A}$ gives
\begin{equation}
    \hat{A} \approx A \cdot N_0^{-\alpha} \cdot (N_0+c_a)^{\hat{\alpha}} = A \left( \frac{N_0+c_a}{N_0} \right)^{\hat{\alpha}} \left(N_0\right)^{\hat{\alpha}-\alpha}.
\end{equation}
Assuming for simplicity that the pivot point is chosen such that the $N_0^{\hat{\alpha}-\alpha}$ term is of order one, we focus on the dominant term
\begin{equation}
    \hat{A} \approx A \left( 1 + \frac{c_a}{N_0} \right)^{\hat{\alpha}}.
\end{equation}
This relationship explains the rapid growth of $\hat{A}$. Since we have already established that $\hat{\alpha}$ itself increases with $c$, the prefactor $\hat{A}$ grows due to two compounding effects: an increase in the base $(1+c_a/N_0)$ and an increase in the exponent $\hat{\alpha}$. This leads to the exponential-like growth observed empirically in~\cref{fig:perturbed_fits}.



\subsubsection{Systematic Bias Perturbation}
\label{app:theoretical_analysis:subsec:systematic}

Assume a bias where the error itself scales with model size, modeled as $\tilde{N} = \mu_{\text{geo}} (N/\mu_{\text{geo}})^s$, where $\mu_{\text{geo}}$ is the geometric mean of the true parameter counts and $s$ is the bias factor. The model size term becomes
\begin{equation}
    \hat{A} \tilde{N}^{-\hat{\alpha}} = \hat{A} \left( \mu_{\text{geo}}^{1-s} N^s \right)^{-\hat{\alpha}} = \left(\hat{A} \mu_{\text{geo}}^{-(1-s)\hat{\alpha}}\right) N^{-s\hat{\alpha}}
\end{equation}
To match the true term $A N^{-\alpha}$, the exponent and the constant term must satisfy the relations
\begin{equation}
    \hat{\alpha} = \frac{\alpha}{s},
    \quad
    \hat{A}
    =
    \mu_{\text{geo}} ^{\frac{\alpha  (1-s)}{s}} A .
\end{equation}
The fitted exponent is inversely proportional to the bias factor $s$, which is verified empirically in~\cref{sec:robustness_analysis:subsec:systematic_bias}. The exponent in the optimal ratio is now $(\alpha/s - \beta)/(\alpha/s + \beta)$.

\textbf{Conclusion:} A systematic bias also breaks the $\hat{\alpha} \approx \beta$ condition, unless $s=1$.
\begin{itemize}
    \item If $s < 1$ (inflating larger models relative to smaller ones), then $\hat{\alpha} > \alpha \approx \beta$. The exponent on $C$ becomes positive, and the optimal ratio \textit{increases} with compute.
    \item If $s > 1$ (shrinking larger models relative to smaller ones), then $\hat{\alpha} < \alpha \approx \beta$. The exponent on $C$ becomes negative, and the optimal ratio \emph{decreases} with compute.
\end{itemize}
This perturbation also qualitatively changes the optimal scaling strategy, with the direction of the change depending on the nature of the bias, as seen in \cref{fig:perturbed_compute_optimal_tpp} (bottom left).

\clearpage